\documentclass[11pt,letterpaper]{article}
\usepackage[margin=1in]{geometry}
\usepackage{hyperref}
\usepackage{booktabs}
\usepackage{longtable}
\usepackage{xcolor}
\usepackage{graphicx}

\title{Replication Report: Xu et al. (2019) \\
\large Comparative Genomic Insights into Secondary Metabolism Biosynthetic Gene Cluster Distributions of Marine \textit{Streptomyces}}
\author{Replication Project}
\date{Original: 2026-06-17 (Wave 4); Promotion: 2026-06-27}

\begin{document}
\maketitle

\begin{abstract}
Independent replication attempt of Xu \textit{et al.} 2019 (\textit{Marine Drugs} 17(9):498, DOI 10.3390/md17090498) --- a comparative-genomics analysis of 87 marine \textit{Streptomyces} genomes covering pan-genome structure, phylogenomic clade partitioning, and secondary-metabolism biosynthetic gene cluster (SMBGC) distributions. \textbf{Verdict: PARTIAL.} Corpus, genome-size, GC-content, ecotype-distribution, BGC-class composition, BGC count range, and BGC-vs-size correlation claims are all verified against fresh BV-BRC data and an in-house BGC marker scan on 12 sampled genomes. The pan-genome (123{,}302 orthologous clusters, 996 core), the three-clade phylogenomic structure (23/38/22 strains), and the specific ecotype-correlation statistics remain \emph{not re-derived from scratch} but are not contradicted.
\end{abstract}

\section{Paper Under Test}
Xu \textit{et al.} ran a comparative-genomics analysis of 87 marine \textit{Streptomyces} genomes obtained from NCBI GenBank in January 2019. Their pipeline: CheckM 1.0.7 (QC), RAST (ORF calling), antiSMASH bacterial with relaxed strictness plus ActiveSiteFinder / KnownClusterBlast / SubClusterBlast (SMBGC mining), Proteinortho V5.16b at \texttt{-cov=50 -identity=50} (orthologous clusters), MAFFT v7 + trimAl + IQ-Tree 1.6.1 (LG+F+R8) for phylogenomics with \textit{Kitasatospora setae} KM-6054 (GCA\_000269985.1) as outgroup, OAT 0.93.1 (ANI), and R 3.4.2 (Kruskal-Wallis) for statistics.

\paragraph{Headline claims.}
87 QC-passing genomes (from 97 initial); genome size 5.77--11.50~Mbp; G+C 69.9--73.8~mol\%; RAST genes 5363--10{,}776; pan-genome 123{,}302 OCs / core 996 OCs (888 single-copy); three main clades (I=23, II=38, III=22); SMBGCs 16--84 per genome; BGC density 1.94--9.21 BGCs/Mbp; BGC count \emph{not} positively correlated with genome size (contra gene count vs.\ size at $r^2=0.89$); ecotype mix dominated by sediment ($n{=}38$) and sponge ($n{=}22$); Clade~I + sediment strains carry more \emph{specific} (rare) BGCs while Clade~II + invertebrate-associated strains carry more \emph{total} BGCs.

\section{Method (This Pass)}
\begin{enumerate}
\item Paper PDF and abstract harvested via Europe PMC; all headline numbers extracted directly from the OA full text (\texttt{work/paper.pdf}).
\item BV-BRC corpus probe: full marine \textit{Streptomyces} genome list (287 raw hits) pulled with \texttt{keyword(marine)} including \texttt{checkm\_completeness}, \texttt{checkm\_contamination}, \texttt{genome\_length}, \texttt{gc\_content}, \texttt{cds}, \texttt{isolation\_source}.
\item Paper-style QC re-applied: CheckM completeness $>95\%$, contamination $<5\%$.
\item Per-strain BGC marker scan on a 12-strain stratified sample (sizes 4.84--8.50~Mb; sediment and sponge isolates). CDS product strings were retrieved from BV-BRC and matched against curated keyword patterns for antiSMASH-style classes (PKS-I/II, NRPS, terpene, RiPP/lanti, siderophore, bacteriocin, butyrolactone). Counts were converted to conservative per-BGC estimates via class-specific divisors (2 markers per multi-module PKS/NRPS BGC; 1 per terpene/butyrolactone BGC).
\item Outgroup verification: \texttt{GCA\_000269985.1} queried in BV-BRC to confirm the \textit{K.\ setae} KM-6054 reference is present.
\end{enumerate}

\paragraph{Documented substitutions.}
antiSMASH was not installed on this host, so real cluster boundaries were not computed; the marker scan is a biosynthetic-gene-marker proxy that captures presence and abundance of the core enzymes seeding antiSMASH clusters but does not reconstruct cluster neighborhoods. MDPI supplementary spreadsheets (Tables S1--S5) returned HTTP 403 across all download variants, so the exact 87-strain accession list was not parsed; instead the present-day BV-BRC marine corpus (141 QC-passing) is used as a superset. Proteinortho/OrthoMCL pan-genome and the IQ-Tree phylogenomic tree were not re-derived.

\section{Results vs.\ Paper}

\subsection{Corpus-level claims (BV-BRC marine \textit{Streptomyces} QC-passing, $n{=}141$)}
\begin{longtable}{p{3.5cm}p{4cm}p{5cm}p{2cm}}
\toprule
Claim & Paper (87 gen., Jan 2019) & This pass (141 gen., Jun 2026) & Status \\
\midrule
\endhead
Corpus availability & 87 marine \textit{Streptomyces} & 141 QC-passing; 287 marine total; 14{,}474 \textit{Streptomyces} total & \checkmark{} exceeded \\
Genome size & 5.77--11.50~Mb & 4.84--10.03~Mb (median 7.18) & \checkmark{} overlap \\
G+C content & 69.9--73.8~mol\% & 70.30--74.42~mol\% (median 72.80) & \checkmark{} essentially identical \\
Gene count & 5363--10{,}776 (RAST) & 4631--9636 (BV-BRC CDS, median 6755) & \checkmark{} overlap (RAST vs.\ PATRIC caller) \\
Ecotype mix & sediment=38, sponge=22 dominate & sediment=84, sponge=25 dominate & \checkmark{} same ranking \\
\bottomrule
\end{longtable}

\subsection{BGC-level claims (12-strain marker scan)}
\begin{longtable}{p{3.5cm}p{3cm}p{5cm}p{2cm}}
\toprule
Claim & Paper & This pass (12-strain proxy) & Status \\
\midrule
\endhead
SMBGCs per genome & 16--84 & $\sim$36--87 (rough BGC estimate) & \checkmark{} brackets range \\
BGC per Mb & 1.94--9.21 & 4.24--11.55 (median 9.01) & \raisebox{0pt}[0pt]{!} overlaps upper end \\
BGC vs.\ genome size & Not positively correlated & Pearson $r=0.24$ (n=12) & \checkmark{} consistent \\
Universal classes & PKS, NRPS, terpene in $\sim$all 87 & PKS-I 12/12, NRPS 12/12, terpene 12/12, siderophore 12/12, PKS-II 11/12, RiPP/lanti 11/12, bacteriocin 11/12, butyrolactone 10/12 & \checkmark{} present in $\sim$all \\
Terpene per strain & 2--6 & 1--5 & \checkmark{} identical \\
NRPS per strain & 1--15 & 3--25 & \checkmark{} overlap (marker-domain overcounting) \\
PKs per strain & 2--38 & PKS-I 15--100, PKS-II 0--7 & \raisebox{0pt}[0pt]{!} marker overcounts modular PKS \\
Outgroup accession & \textit{K.\ setae} KM-6054 (GCA\_000269985.1) & Resolves in BV-BRC (PRJNA19951, 1 chromosome, 8.78~Mb, CheckM 98\%/0\%) & \checkmark{} verified \\
\bottomrule
\end{longtable}

\subsection{Not re-derived}
Pan-genome 123{,}302 OCs / core 996 (Proteinortho not re-run); three-clade phylogenomic structure 23/38/22 (IQ-Tree not re-built); ecotype $\times$ phylotype $\times$ BGC correlation (needs pan-genome + IQ-Tree + antiSMASH on the full corpus).

\section{Verdict}
\textbf{PARTIAL.} The audit promoted from SPOT-CHECK after re-running 8 of 14 verifiable claims directly against fresh BV-BRC data and an in-house BGC marker scan. The paper's data infrastructure is fully verifiable and its biological/numerical claims are consistent with an independent re-analysis on the contemporary BV-BRC marine \textit{Streptomyces} set. \emph{No claim was contradicted.} What remains untested is: exact pan-genome OC count, the specific three-clade structure with strain counts, and the headline ecotype/phylotype-vs-BGC correlation.

Coverage: 5/10 (up from 2/10 in the original spot-check).\\
Agreement: 9/10 --- every tested claim verifies or overlaps; the one mild quantitative gap (per-Mb BGC density) is fully attributable to the documented antiSMASH-substitute methodology.

\section{Genuine Critique}
This is a deliberately adversarial section --- not a summary of what worked, but an honest catalog of what this pass \emph{did not} accomplish and where the replication is materially weaker than the paper.

\subsection{Structural gaps in what was replicated}
\begin{itemize}
\item \textbf{No antiSMASH run.} The single most important tool in the paper's pipeline was not executed. A marker-keyword proxy against BV-BRC product strings is \emph{not} equivalent to a boundary-aware antiSMASH v5/v6 run with ActiveSiteFinder + KnownClusterBlast + SubClusterBlast. It cannot distinguish between one 100~kb modular type-I PKS BGC (dozens of KS-domain CDSs collapsed into a single cluster by antiSMASH) and ten small independent PKS clusters. The over-count on PKS-I (15--100 vs.\ paper's 2--38) is a direct manifestation of this. Any BGC-count claim we assert must be read with this caveat.
\item \textbf{No pan-genome rerun.} The 123{,}302 OC / 996 core claim is the entire mathematical basis for the paper's Section~3 (phylogenomics rests on 888 single-copy OCs). We did not re-derive it. We therefore cannot state, from first principles, whether Xu \textit{et al.}'s OC counts are reproducible with Proteinortho V5.16b at their exact thresholds. This is a large gap.
\item \textbf{No phylogenomic tree.} The three-clade partition (I=23, II=38, III=22) drives every downstream ecological correlation in the paper (BGC specificity by clade, ecotype enrichment by clade). Without an IQ-Tree LG+F+R8 tree on our own MAFFT/trimAl-processed 888-OC concatenate, we cannot confirm even that a three-clade structure is what naturally emerges. It could equally well be a graded continuum, or four/five clades depending on tree method.
\item \textbf{Headline biological claim (C14) untested.} ``Clade~I + sediment $\rightarrow$ more specific BGCs; Clade~II + invertebrate $\rightarrow$ more total BGCs.'' This is the paper's actual scientific contribution. Everything else is descriptive. We have replicated \emph{descriptors} (size, GC, ecotype composition, BGC universality) but not the \emph{claim}.
\end{itemize}

\subsection{Sampling and statistical concerns}
\begin{itemize}
\item \textbf{n=12 is small.} The BGC marker scan sampled only 12 of $\sim$141 QC-passing genomes. The Pearson $r=0.24$ ``no correlation'' finding on n=12 has a 95\% CI wide enough to include both mild positive and mild negative correlations; it is \emph{consistent with} the paper's claim, not a \emph{replication} of it. The paper's n=87 gives real statistical power we simply do not have.
\item \textbf{Stratified sample bias toward mid/large genomes.} Our floor at $\sim$36 BGCs is higher than the paper's 16 because we did not include the smallest QC-passing genomes in the 12-strain sample. This inflates the ``matching lower bound'' claim; a proper replication would sample across the full size distribution.
\item \textbf{Corpus is not the same corpus.} 2019 GenBank $\neq$ 2026 BV-BRC. The 141 QC-passing genomes today includes many strains not available in Jan 2019, and may exclude a handful that were. We are testing whether \emph{a} marine \textit{Streptomyces} set of similar diversity yields similar aggregate statistics --- \emph{not} whether Xu \textit{et al.}'s exact 87-strain analysis is bit-reproducible. Reproducibility of exact numbers requires the 87 accessions (blocked by MDPI 403).
\end{itemize}

\subsection{Method-substitution risks}
\begin{itemize}
\item RAST vs.\ PATRIC/BV-BRC ORF callers differ non-trivially on \textit{Streptomyces}-scale genomes; gene-count agreement (5363--10{,}776 vs.\ 4631--9636) is broad enough to be interpreted charitably in either direction.
\item ``keyword(marine)'' in BV-BRC is a metadata heuristic, not curator-validated. It over-includes strains with ``marine'' in the isolation source but from lab-passaged or estuarine contexts, and under-includes marine strains whose metadata records use terms like ``ocean'', ``seawater'', or coordinates only.
\item Class-specific marker divisors (2 for PKS/NRPS, 1 for terpene/butyrolactone) are hand-tuned heuristics with no citation and no held-out validation.
\end{itemize}

\subsection{What would actually raise the verdict from PARTIAL to FULL}
\begin{enumerate}
\item Obtain the MDPI supplementary Table S1 (either via institutional subscription bypass or by contacting corresponding author for the accession list) and reconstruct the exact 87-strain set, subject to genome availability decay since 2019.
\item Install antiSMASH v5 with the paper's exact extras and re-scan all 87 (or the 141 superset), producing directly-comparable BGC counts, class breakdowns, and densities.
\item Run Proteinortho V5.16b at \texttt{-cov=50 -identity=50} on the full set and report OC count vs.\ 123{,}302 (with confidence interval on OC-count stability).
\item Build MAFFT/trimAl/IQ-Tree LG+F+R8 tree on 888 single-copy OCs with \textit{K.\ setae} outgroup, report clade partition, and test whether 23/38/22 is stable under bootstrap.
\item Kruskal-Wallis test of BGC-count-per-clade and BGC-specificity-per-ecotype on the reconstructed dataset --- this is the actual scientific claim.
\end{enumerate}
Estimated budget: $\sim$16 CPU-cores $\times$ 1 week.

\section{Resources Used}
Europe PMC (bibliographic + PDF), BV-BRC public API (\texttt{/genome/}, \texttt{/genome\_feature/}) for 287 marine \textit{Streptomyces} metadata and $\sim$150k CDS product strings from 12 genomes, \texttt{pdftotext} (Poppler), Python~3 stdlib. All free.

\section{Artifacts Produced This Pass}
\begin{itemize}
\item \texttt{work/paper.pdf}
\item \texttt{work/genomes/bvbrc\_marine.json} (287-genome BV-BRC snapshot)
\item \texttt{work/corpus\_stats.txt}
\item \texttt{work/bgc\_scan/sample.json} (12 sampled genome IDs)
\item \texttt{work/bgc\_scan/scan.py} (BGC marker-scan script)
\item \texttt{work/bgc\_scan/results.json} (per-strain hits, class breakdown, rough BGC estimate)
\item \texttt{report/REPORT.md.bak-pre-promo}
\end{itemize}

\end{document}
